\documentclass[11pt,a4paper]{article}

% ─── Packages ───────────────────────────────────────────────
\usepackage[utf8]{inputenc}
\usepackage[T1]{fontenc}
\usepackage{amsmath,amssymb,amsthm}
\usepackage{mathtools}
\usepackage{algorithm}
\usepackage{algpseudocode}
\usepackage{graphicx}
\usepackage{hyperref}
\usepackage[margin=2.5cm]{geometry}
\usepackage{booktabs}
\usepackage{caption}
\usepackage{subcaption}
\usepackage{enumitem}

% ─── Theorem environments ───────────────────────────────────
\newtheorem{theorem}{Theorem}
\newtheorem{lemma}[theorem]{Lemma}
\newtheorem{proposition}[theorem]{Proposition}
\newtheorem{remark}{Remark}

% ─── Custom commands ────────────────────────────────────────
\newcommand{\R}{\mathbb{R}}
\newcommand{\norm}[1]{\left\|#1\right\|}
\newcommand{\inner}[2]{\langle #1, #2 \rangle}
\newcommand{\Xhat}{\hat{X}}
\newcommand{\yhat}{\hat{y}}
\DeclareMathOperator{\diag}{diag}
\DeclareMathOperator{\sign}{sign}

% ═══════════════════════════════════════════════════════════════
\begin{document}

% ─── Title ──────────────────────────────────────────────────
\begin{center}
    {\Large\bfseries Computational Mathematics for Learning and Data Analysis}\\[6pt]
    {\large Project Track 19 --- Non-ML}\\[12pt]
    {\large\bfseries QR Factorization and Limited-Memory \\Quasi-Newton Methods for Regularized Least Squares}\\[12pt]
    
    \textit{Karol Jinet}\\
    \texttt{k.jinet@studenti.unipi.it}\\[6pt]
    Academic Year 2024--2025
\end{center}

\vspace{1cm}

% ─── Table of Contents ──────────────────────────────────────
\tableofcontents
\newpage

% ═════════════════════════════════════════════════════════════
\section{Introduction}
\label{sec:intro}
% ═════════════════════════════════════════════════════════════

This project addresses the numerical solution of a regularized linear least squares problem. 
Given a data matrix $X \in \R^{m \times n}$ from the ML-CUP dataset, a random response vector 
$y \in \R^n$, and a regularization parameter $\lambda > 0$, the goal is to find the vector 
$w \in \R^m$ that minimizes

\begin{equation}
\label{eq:problem}
    \min_w \norm{\Xhat w - \yhat},
\end{equation}

where the augmented matrix and vector are defined as

\begin{equation}
\label{eq:augmented}
    \Xhat = \begin{bmatrix} X^T \\ \lambda I_m \end{bmatrix} \in \R^{(n+m) \times m}, 
    \qquad
    \yhat = \begin{bmatrix} y \\ 0 \end{bmatrix} \in \R^{n+m}.
\end{equation}

Expanding the squared norm, the problem is equivalent to minimizing

\begin{equation}
\label{eq:expanded}
    f(w) = \frac{1}{2}\norm{X^T w - y}^2 + \frac{1}{2}\lambda^2 \norm{w}^2,
\end{equation}

which is a standard instance of Tikhonov regularization (ridge regression)~\cite{hastie2009elements}. 
The regularization term $\lambda^2\norm{w}^2$ penalizes large coefficient vectors, 
improving the numerical stability of the solution and controlling overfitting.

Two computational strategies are explored:

\begin{itemize}[nosep]
    \item \textbf{(A1)} A limited-memory quasi-Newton method (L-BFGS), which iteratively 
    approximates the curvature of the objective function using only gradient information 
    and a small number of stored vector pairs.
    \item \textbf{(A2)} A direct method based on thin QR factorization with Householder 
    reflectors, in the compact variant where the orthogonal matrix~$Q$ is not formed 
    explicitly, but applied implicitly through stored Householder vectors.
\end{itemize}

Both algorithms are implemented from scratch, without relying on any off-the-shelf 
numerical solvers. The first part of the report studies the iterative method (A1), 
while the second part derives the direct QR-based solver (A2).


% ═════════════════════════════════════════════════════════════
\section{Problem Properties}
\label{sec:properties}
% ═════════════════════════════════════════════════════════════

Before discussing the optimization algorithm, we analyze the mathematical properties 
of the objective function~\eqref{eq:expanded} that are relevant for convergence.

\subsection{Gradient and Hessian}

The objective~\eqref{eq:expanded} is a quadratic function with constant Hessian. 
Its gradient and Hessian are:

\begin{align}
    \nabla f(w) &= X(X^T w - y) + \lambda^2 w, \label{eq:gradient} \\
    H \coloneqq \nabla^2 f(w) &= XX^T + \lambda^2 I_m. \label{eq:hessian}
\end{align}

\subsection{Strong Convexity and Smoothness}

Since $XX^T \succeq 0$ and $\lambda > 0$, the Hessian is symmetric positive definite:

\begin{equation}
    v^T H v = \norm{X^T v}^2 + \lambda^2\norm{v}^2 \geq \lambda^2\norm{v}^2 > 0,
    \quad \forall\, v \neq 0.
\end{equation}

This implies that $f$ is $\mu$-strongly convex with parameter $\mu = \lambda_{\min}(H) = \lambda_m(XX^T) + \lambda^2$, 
where $\lambda_m(XX^T)$ denotes the smallest eigenvalue of $XX^T$. Since $X$ is tall-thin 
($m > n$), the matrix $XX^T$ is rank-deficient with $\lambda_m(XX^T) = 0$ 
(see Lemma~\ref{lemma:eigenvalues} in Appendix), hence $\mu = \lambda^2$.

The gradient is Lipschitz continuous with constant $L = \lambda_{\max}(H) = \lambda_1(XX^T) + \lambda^2$, 
since for all $w, w' \in \R^m$:

\begin{equation}
    \norm{\nabla f(w) - \nabla f(w')} = \norm{H(w - w')} \leq \norm{H}\norm{w - w'} = L\norm{w - w'}.
\end{equation}

\subsection{Condition Number}
\label{sec:conditioning}

The condition number of the augmented matrix $\Xhat$ is

\begin{equation}
\label{eq:condition}
    \kappa(\Xhat) = \frac{\sigma_{\max}(\Xhat)}{\sigma_{\min}(\Xhat)} 
    = \sqrt{\frac{\lambda_1(XX^T) + \lambda^2}{\lambda^2}} 
    = \frac{\sqrt{\lambda_1(XX^T) + \lambda^2}}{\lambda},
\end{equation}

which scales as $O(1/\lambda)$ for small $\lambda$. This means that weak regularization 
leads to ill-conditioned problems, directly affecting the convergence speed of iterative methods.


% ═════════════════════════════════════════════════════════════
\section{Limited-Memory BFGS (L-BFGS)}
\label{sec:lbfgs}
% ═════════════════════════════════════════════════════════════

\subsection{Quasi-Newton Methods: Background}

Quasi-Newton methods are iterative algorithms for unconstrained optimization that 
approximate the Newton direction without explicitly computing the Hessian 
matrix~\cite{nocedal2006numerical}. At each iteration~$k$, the search direction is 
computed as

\begin{equation}
\label{eq:qn_direction}
    p_k = -H_k \nabla f(w_k),
\end{equation}

where $H_k \approx [\nabla^2 f(w_k)]^{-1}$ is an approximation to the inverse Hessian, 
built incrementally from gradient differences. The iterate is then updated as

\begin{equation}
    w_{k+1} = w_k + \alpha_k p_k,
\end{equation}

where $\alpha_k > 0$ is a step length determined by a line search procedure.

The BFGS update~\cite{broyden1970} maintains a dense $m \times m$ matrix $H_k$ 
that is updated at each iteration via a rank-two formula using the pairs

\begin{equation}
    s_k = w_{k+1} - w_k, \qquad y_k = \nabla f(w_{k+1}) - \nabla f(w_k).
\end{equation}

While BFGS achieves superlinear convergence on strongly convex 
problems~\cite{nocedal2006numerical}, storing and updating the full $m \times m$ 
matrix becomes prohibitive when the dimension~$m$ is large.


\subsection{The L-BFGS Algorithm}

The Limited-Memory BFGS (L-BFGS) method~\cite{liu1989limited} addresses the memory 
limitation by storing only the most recent $\bar{m}$ pairs $\{s_i, y_i\}$, typically 
with $\bar{m} \in [3, 20]$. Rather than forming the full matrix $H_k$, the product 
$H_k \nabla f(w_k)$ is computed implicitly via a \emph{two-loop recursion} 
(Algorithm~\ref{alg:twoloop}).

\subsubsection{Initial Scaling}

At each iteration, the initial approximation $H_k^0 = \gamma_k I$ is set 
using~\cite[Eq.~9.6]{nocedal2006numerical}:

\begin{equation}
\label{eq:gamma}
    \gamma_k = \frac{s_{k-1}^T y_{k-1}}{y_{k-1}^T y_{k-1}}.
\end{equation}

This scaling attempts to estimate the magnitude of the true Hessian along the most 
recent search direction, ensuring that the search direction $p_k$ is well-scaled 
and that the unit step length $\alpha_k = 1$ is accepted in most iterations.


\subsubsection{Two-Loop Recursion}

The product $r = H_k \nabla f_k$ is computed efficiently by the following 
procedure~\cite[Algorithm~9.1]{nocedal2006numerical}:

\begin{algorithm}[H]
\caption{L-BFGS Two-Loop Recursion~\cite[Algorithm~9.1]{nocedal2006numerical}}
\label{alg:twoloop}
\begin{algorithmic}[1]
    \Require Gradient $q \gets \nabla f_k$; stored pairs $\{s_i, y_i\}_{i=k-\bar{m}}^{k-1}$
    \Ensure $r = H_k \nabla f_k$
    \For{$i = k-1, k-2, \ldots, k-\bar{m}$}
        \State $\rho_i \gets 1 / (y_i^T s_i)$
        \State $\alpha_i \gets \rho_i\, s_i^T q$
        \State $q \gets q - \alpha_i\, y_i$
    \EndFor
    \State $r \gets H_k^0\, q$ \Comment{$H_k^0 = \gamma_k I$ from Eq.~\eqref{eq:gamma}}
    \For{$i = k-\bar{m}, \ldots, k-1$}
        \State $\beta \gets \rho_i\, y_i^T r$
        \State $r \gets r + s_i\,(\alpha_i - \beta)$
    \EndFor
    \State \Return $r$
\end{algorithmic}
\end{algorithm}

The cost of the two-loop recursion is $O(4\bar{m}\,m)$ multiplications, plus $O(m)$ for 
the scaling $H_k^0 q$. This is a significant improvement over the $O(m^2)$ cost of 
full BFGS.


\subsubsection{Curvature Safeguard}

The BFGS update requires that the curvature condition $y_k^T s_k > 0$ holds. 
When this condition is satisfied, the pair $(s_k, y_k)$ is added to the memory; 
otherwise, it is discarded. If the curvature becomes excessively small 
($y_k^T s_k \leq 10^{-16}$), the stored memory is cleared entirely and the algorithm 
restarts from a gradient descent step, following~\cite{liu1989limited}.


% ─────────────────────────────────────────────────────────────
\subsection{Line Search Strategies}
\label{sec:linesearch}
% ─────────────────────────────────────────────────────────────

\subsubsection{Strong Wolfe Conditions}

The standard approach for selecting the step length $\alpha_k$ in quasi-Newton methods 
is a line search satisfying the \emph{strong Wolfe conditions}~\cite[Eq.~3.7]{nocedal2006numerical}:

\begin{align}
    f(w_k + \alpha_k p_k) &\leq f(w_k) + c_1\, \alpha_k\, \nabla f_k^T p_k, \label{eq:armijo} \\
    \left|\nabla f(w_k + \alpha_k p_k)^T p_k\right| &\leq c_2\, \left|\nabla f_k^T p_k\right|, \label{eq:curvature}
\end{align}

with $0 < c_1 < c_2 < 1$. Following standard practice for quasi-Newton 
methods~\cite{nocedal2006numerical}, we set $c_1 = 10^{-4}$ and $c_2 = 0.9$.

Condition~\eqref{eq:armijo} (Armijo) ensures sufficient decrease in the objective value, 
while~\eqref{eq:curvature} prevents the step from being too short by requiring that the 
directional derivative at the new point is sufficiently small.

Our implementation uses a bracketing-and-zoom procedure with cubic interpolation, 
following the algorithmic framework described in~\cite[Section~3.5]{nocedal2006numerical}.


\subsubsection{Exact Line Search for Quadratic Objectives}
\label{sec:exact_ls}

Since the objective function~\eqref{eq:expanded} is quadratic, the restriction of $f$ 
along any direction $p$ is a parabola:

\begin{equation}
    \varphi(\alpha) = f(w + \alpha\, p) 
    = f(w) + \alpha\,\nabla f(w)^T p + \frac{\alpha^2}{2}\,p^T H p.
\end{equation}

The second derivative $\varphi''(\alpha) = p^T H p$ is constant and positive (since $H \succ 0$), 
so $\varphi$ has a unique global minimum at

\begin{equation}
\label{eq:exact_alpha}
    \alpha_{\min} = -\frac{\nabla f(w)^T p}{p^T H p} 
    = -\frac{\nabla f(w)^T p}{\norm{X^T p}^2 + \lambda^2\norm{p}^2}.
\end{equation}

This is a generalization of~\cite[Eq.~3.25]{nocedal2006numerical} to an arbitrary descent 
direction~$p$. The computational cost is dominated by a single matrix-vector product 
$X^T p$, requiring $O(mn)$ operations --- the same cost as a gradient evaluation, and 
considerably cheaper than the multiple function and gradient evaluations typically needed 
by the Wolfe line search.

For quadratic problems, the exact line search is theoretically preferable because it 
minimizes $f$ exactly along each search direction, fully exploiting the parabolic structure 
of the objective. As noted by~\cite{liu1989limited} and confirmed by our experiments 
(Section~\ref{sec:exact_vs_wolfe}), this leads to faster convergence and higher final accuracy.


% ─────────────────────────────────────────────────────────────
\subsection{Complete Algorithm}
% ─────────────────────────────────────────────────────────────

The complete L-BFGS procedure is summarized in Algorithm~\ref{alg:lbfgs}.

\begin{algorithm}[H]
\caption{L-BFGS~\cite[Algorithm~9.2]{nocedal2006numerical}}
\label{alg:lbfgs}
\begin{algorithmic}[1]
    \Require Starting point $w_0$; memory size $\bar{m} > 0$; tolerance $\varepsilon > 0$
    \State $k \gets 0$; compute $\nabla f_0$; set $\texttt{stop\_tol} \gets \varepsilon\,\norm{\nabla f_0}$
    \While{$\norm{\nabla f_k} > \texttt{stop\_tol}$}
        \State Choose $H_k^0 = \gamma_k I$ \Comment{Eq.~\eqref{eq:gamma}}
        \State Compute $p_k \gets -H_k \nabla f_k$ \Comment{Algorithm~\ref{alg:twoloop}}
        \If{$\nabla f_k^T p_k \geq 0$} \Comment{Safeguard: fallback to steepest descent}
            \State $p_k \gets -\nabla f_k$
        \EndIf
        \State Compute $\alpha_k$ via exact line search (Eq.~\eqref{eq:exact_alpha}) 
               or Wolfe conditions
        \State $w_{k+1} \gets w_k + \alpha_k\, p_k$
        \State $s_k \gets w_{k+1} - w_k$; \quad $y_k \gets \nabla f_{k+1} - \nabla f_k$
        \If{$y_k^T s_k > 10^{-10}$}
            \State Store $(s_k, y_k)$; discard oldest pair if $\#\text{stored} > \bar{m}$
        \ElsIf{$y_k^T s_k \leq 10^{-16}$}
            \State Clear all stored pairs \Comment{Curvature reset}
        \EndIf
        \State $k \gets k + 1$
    \EndWhile
    \State \Return $w_k$
\end{algorithmic}
\end{algorithm}


% ─────────────────────────────────────────────────────────────
\subsection{Stopping Criterion}
\label{sec:stopping}
% ─────────────────────────────────────────────────────────────

The primary stopping criterion uses a \emph{relative} tolerance on the gradient norm:

\begin{equation}
\label{eq:stop_rel}
    \norm{\nabla f(w_k)} \leq \varepsilon \cdot \norm{\nabla f(w_0)},
\end{equation}

where $\varepsilon > 0$ is a user-specified tolerance and $w_0$ is the starting point. 
This criterion is scale-invariant: it does not depend on the absolute magnitude of the 
gradient, which can vary significantly across different values of~$\lambda$.

As a secondary criterion, the algorithm detects \emph{numerical stagnation}: 
if the relative change in function value $|f_{k} - f_{k-1}|/|f_k|$ drops below 
machine epsilon while the step length $\alpha_k$ is near zero, the iteration is terminated.


% ─────────────────────────────────────────────────────────────
\subsection{Convergence Analysis}
\label{sec:convergence}
% ─────────────────────────────────────────────────────────────

\subsubsection{Global Linear Convergence}

The convergence of L-BFGS on the problem at hand follows from 
\cite[Theorem~1]{liu1989limited}. The key assumptions are:

\begin{enumerate}[nosep]
    \item $f \in C^2$ (satisfied: $f$ is quadratic);
    \item the level sets $\mathcal{L} = \{w : f(w) \leq f(w_0)\}$ are convex 
          (satisfied: $f$ is strongly convex);
    \item $\exists\, M_1, M_2 > 0$ such that $M_1\norm{v}^2 \leq v^T \nabla^2 f(w)\,v \leq M_2\norm{v}^2$ 
          for all $v \in \R^m$ and $w \in \mathcal{L}$ 
          (satisfied: $H$ is constant with $M_1 = \lambda^2$ and $M_2 = \lambda_1(XX^T) + \lambda^2$).
\end{enumerate}

Under these conditions, the iterates converge linearly:

\begin{equation}
    f(w_k) - f(w^*) \leq r^k\,(f(w_0) - f(w^*)), \quad r \in [0,1).
\end{equation}

\subsubsection{Superlinear Convergence}

For the full BFGS method (i.e., when $\bar{m}$ is large enough to store the entire history), 
\cite[Theorem~6.6]{nocedal2006numerical} guarantees superlinear convergence on strongly 
convex functions.

A non-asymptotic analysis by Rodomanov and Nesterov~\cite{rodomanov2021} predicts that 
superlinear convergence begins after approximately $K_0^{\text{BFGS}} = 4m\ln\kappa$ iterations, 
where $\kappa = L/\mu$ is the condition number of the Hessian. 
Jin and Mokhtari~\cite{jin2022} provide a local bound:

\begin{equation}
    \frac{\norm{w_k - w^*}_H}{\norm{w_0 - w^*}_H} \leq \left(\frac{1}{k}\right)^{k/2},
\end{equation}

where $\norm{v}_H = \sqrt{v^T H v}$ is the Hessian-induced norm. This implies that 
achieving accuracy $\varepsilon$ requires only 
$k = O\!\left(\frac{\log(1/\varepsilon)}{\log\log(1/\varepsilon)}\right)$ iterations, 
which is significantly better than the $O(\log(1/\varepsilon))$ bound for linear convergence.

\begin{remark}
For L-BFGS with finite memory $\bar{m}$, the theoretical guarantee is only linear convergence. 
However, in practice, superlinear behavior is often observed for well-conditioned problems, 
as the limited-memory approximation captures sufficient curvature information when $\bar{m}$ 
is large enough relative to the problem dimension.
\end{remark}

\subsubsection{Computational Complexity}

Each iteration of L-BFGS requires:
\begin{itemize}[nosep]
    \item One gradient evaluation: $O(mn)$ (matrix-vector product $X(X^T w - y)$);
    \item Two-loop recursion: $O(4\bar{m}\,m)$;
    \item Line search: $O(mn)$ for exact (one product $X^T p$), or $O(k_{LS}\,mn)$ for Wolfe 
          (where $k_{LS}$ is the number of line search evaluations).
\end{itemize}

The total cost per iteration is $O(mn + \bar{m}\,m)$, and the memory requirement is 
$O(\bar{m}\,m)$ for storing the vector pairs --- a drastic improvement over the $O(m^2)$ 
storage of full BFGS.


% ═════════════════════════════════════════════════════════════
\section{L-BFGS Experimental Results}
\label{sec:experiments}
% ═════════════════════════════════════════════════════════════

All experiments use a matrix $X \in \R^{m \times n}$ generated randomly with i.i.d.\ standard 
normal entries (as a proxy for the ML-CUP dataset), with default dimensions $m = 600$ 
and $n = 12$. The response vector $y \in \R^n$ is drawn from a standard normal distribution 
with a fixed random seed. Unless stated otherwise, the tolerance is $\varepsilon = 10^{-14}$ 
(relative), the memory size is $\bar{m} = 10$, and the baseline solution $w^*$ is computed 
via the normal equations using \texttt{numpy.linalg.solve}.


\subsection{Exact Line Search vs.\ Strong Wolfe Line Search}
\label{sec:exact_vs_wolfe}

We compare the two line search strategies on the same problem instance with $\lambda = 0.5$.

\begin{table}[h]
\centering
\caption{Comparison of line search strategies ($m=600$, $n=12$, $\lambda=0.5$, $\varepsilon=10^{-14}$ relative).}
\label{tab:ls_comparison}
\begin{tabular}{lcccr}
    \toprule
    \textbf{Line Search} & \textbf{Iterations} & $\norm{w - w^*}$ & $|f^* - f|$ & \textbf{Time (s)} \\
    \midrule
    Exact   & 16 & $6.4 \times 10^{-14}$ & $4.3 \times 10^{-19}$ & 0.002 \\
    Wolfe   & 16 & $6.0 \times 10^{-11}$ & $< 10^{-19}$ & 0.002 \\
    \bottomrule
\end{tabular}
\end{table}

Both strategies converge in the same number of iterations, but the exact line search 
achieves approximately three orders of magnitude higher accuracy on the weight vector 
($10^{-14}$ vs.\ $10^{-11}$). This is because the exact line search optimizes each step 
perfectly along the parabolic profile, while the Wolfe conditions allow a range of 
acceptable step lengths. The Wolfe variant stagnates earlier with a gradient norm 
of $\sim\!4 \times 10^{-8}$, while the exact variant continues to decrease until 
$\sim\!5 \times 10^{-13}$.

%% TODO: insert convergence plot (Figure 1) when notebooks are run


\subsection{Effect of the Memory Parameter \texorpdfstring{$\bar{m}$}{mbar}}
\label{sec:memory_effect}

We vary $\bar{m} \in \{3, 5, 10, 20, 40\}$ using exact line search, relative tolerance 
$\varepsilon = 10^{-12}$, and $\lambda = 0.5$.

For well-conditioned problems (moderate $\lambda$), the effect of $\bar{m}$ is minimal: 
all configurations converge within approximately the same number of iterations. 
This is consistent with the observation by Nocedal and Wright~\cite[Section~9.1]{nocedal2006numerical} 
that modest values of $\bar{m}$ between 3 and 20 often produce satisfactory results.

For ill-conditioned problems (small $\lambda$), larger memory sizes provide a more 
accurate approximation of the curvature, leading to faster convergence.

%% TODO: insert memory comparison plot (Figure 2)


\subsection{Effect of the Regularization Parameter \texorpdfstring{$\lambda$}{lambda}}
\label{sec:lambda_effect}

We study convergence across several regularization regimes with 
$\lambda \in \{10^{-4}, 10^{-2}, 1, 10^2, 10^4\}$, using exact line search and 
$\bar{m} = 10$.

As $\lambda$ increases, the condition number $\kappa(\Xhat)$ decreases 
(cf.\ Eq.~\eqref{eq:condition}), leading to faster convergence. For $\lambda = 10^4$, 
the Hessian is nearly $\lambda^2 I$ and the algorithm converges in fewer than 10~iterations. 
For $\lambda = 10^{-4}$, the condition number exceeds $10^5$, and convergence requires 
significantly more iterations.

This behavior empirically validates the theoretical dependence of the convergence rate 
on the condition number $\kappa$, as predicted by~\cite{rodomanov2021} and~\cite{jin2022}.

%% TODO: insert lambda comparison plot and table (Figure 3, Table 2)


\subsection{Relative Error Trajectory}

To monitor convergence quality, we track the relative error

\begin{equation}
    e_k = \frac{\norm{w_k - w^*}}{\norm{w^*}}
\end{equation}

across iterations for different values of $\lambda$. For well-conditioned problems 
($\lambda \geq 1$), the relative error decays smoothly to machine precision. 
For ill-conditioned problems ($\lambda \leq 10^{-2}$), the decay is slower and may 
exhibit non-monotonic behavior, reflecting the difficulty of accumulating accurate 
curvature information in flat regions of the objective landscape.


% ═════════════════════════════════════════════════════════════
\section{QR Factorization}
\label{sec:qr}
% ═════════════════════════════════════════════════════════════

We now turn from the optimization point of view to the direct linear algebra point 
of view. The same regularized least squares problem can be solved without generating 
an iterative sequence of approximations: one factorizes the augmented matrix

\begin{equation}
    A \coloneqq \Xhat =
    \begin{bmatrix}
        X^T \\ \lambda I_m
    \end{bmatrix}
    \in \R^{(n+m)\times m},
    \qquad
    b \coloneqq \yhat =
    \begin{bmatrix}
        y \\ 0
    \end{bmatrix}
    \in \R^{n+m},
\end{equation}

and then solves the resulting triangular problem. Since $\lambda > 0$, 
Lemma~\ref{lemma:fullrank} shows that $A$ has full column rank. Therefore the 
least squares problem

\begin{equation}
    \min_{w\in\R^m} \norm{Aw-b}
\end{equation}

has a unique minimizer.

\subsection{Thin QR Factorization}

For a full-column-rank matrix $A\in\R^{(n+m)\times m}$, the thin QR 
factorization is

\begin{equation}
\label{eq:thin_qr}
    A = Q_1 R,
    \qquad
    Q_1 \in \R^{(n+m)\times m}, \quad Q_1^T Q_1 = I_m,
    \quad
    R \in \R^{m\times m},
\end{equation}

where $R$ is upper triangular and nonsingular. If $Q=[Q_1\ Q_2]$ is an orthogonal 
completion of $Q_1$, then

\begin{equation}
    Q^T A =
    \begin{bmatrix}
        R \\ 0
    \end{bmatrix},
    \qquad
    Q^T b =
    \begin{bmatrix}
        c \\ d
    \end{bmatrix},
    \quad c\in\R^m,\ d\in\R^{n}.
\end{equation}

Orthogonal transformations preserve the Euclidean norm, hence

\begin{equation}
\label{eq:qr_residual_split}
    \norm{Aw-b}^2
    =
    \norm{Q^T(Aw-b)}^2
    =
    \norm{Rw-c}^2 + \norm{d}^2.
\end{equation}

The term $\norm{d}^2$ is independent of $w$, while $R$ is nonsingular. Therefore the 
least squares minimizer is obtained from the triangular system

\begin{equation}
\label{eq:qr_triangular_solve}
    Rw = c = Q_1^T b.
\end{equation}

This is the QR analogue of solving the normal equations. Indeed,

\begin{equation}
\label{eq:qr_normal_equivalence}
    R^T R = A^T A = XX^T + \lambda^2 I_m,
    \qquad
    R^T c = A^T b = Xy.
\end{equation}

Thus QR and the normal equations define the same exact solution. Numerically, however, 
QR is preferable because it avoids explicitly forming and solving with $A^T A$, whose 
condition number is squared:

\begin{equation}
    \kappa(A^T A) = \kappa(A)^2.
\end{equation}

For small $\lambda$, this distinction is important, since Eq.~\eqref{eq:condition} 
shows that $\kappa(A)$ already grows like $O(1/\lambda)$.

\subsection{Algorithmic Structure}

The QR-based solver can be summarized as follows:

\begin{algorithm}[H]
\caption{QR Solver for the Regularized Least Squares Problem}
\label{alg:qr_solver}
\begin{algorithmic}[1]
    \Require Data matrix $X\in\R^{m\times n}$; response $y\in\R^n$; parameter $\lambda>0$
    \State Form the linear operator $A = [X^T;\lambda I_m]$ and vector $b=[y;0]$
    \State Compute a thin QR factorization $A=Q_1R$
    \State Compute $c \gets Q_1^Tb$
    \State Solve $Rw=c$ by back-substitution
    \State \Return $w$
\end{algorithmic}
\end{algorithm}

In practice, the matrix $Q_1$ is not formed explicitly. Forming it would require 
$O((n+m)m)$ storage and applying it as a dense matrix would obscure the sparse and 
structured nature of $A$. Instead, the orthogonal transformations used during the 
factorization are stored in compact form and later applied to $b$.

\subsection{Cost for the Project Matrix}
\label{sec:qr_cost_project}

The matrix in this project is not a generic dense rectangular matrix. It has the block 
form

\begin{equation}
    A =
    \begin{bmatrix}
        X^T \\ \lambda I_m
    \end{bmatrix},
\end{equation}

where $X$ is tall-thin and the ML-CUP data dimension $n$ is small compared with $m$. 
The lower block is diagonal, while all data-dependent information is contained in the 
$n$ rows of $X^T$. A direct dense Householder factorization that ignores this structure 
would cost

\begin{equation}
    O((n+m)m^2),
\end{equation}

which contains a cubic term in $m$. This is not the right implementation model for the 
project requirement that the QR code should be at most quadratic in $m$.

The useful observation is that the row order is irrelevant for the least squares 
solution. If $P$ is any permutation matrix, then

\begin{equation}
    \norm{Aw-b} = \norm{P(Aw-b)}.
\end{equation}

We may therefore work with the permuted system

\begin{equation}
\label{eq:permuted_augmented}
    \tilde{A} =
    \begin{bmatrix}
        \lambda I_m \\ X^T
    \end{bmatrix},
    \qquad
    \tilde{b} =
    \begin{bmatrix}
        0 \\ y
    \end{bmatrix}.
\end{equation}

The first $m$ rows already form the triangular factor $R_0=\lambda I_m$. The $n$ dense 
rows of $X^T$ can then be inserted one at a time by triangularizing

\begin{equation}
    \begin{bmatrix}
        R_{i-1} \\ x_i^T
    \end{bmatrix},
    \qquad i=1,\ldots,n,
\end{equation}

where $x_i$ denotes the $i$-th column of $X$, equivalently the $i$-th row of $X^T$. 
Each row insertion costs $O(m^2)$, so the total cost is

\begin{equation}
\label{eq:structured_qr_cost}
    O(nm^2),
\end{equation}

plus $O(m^2)$ for the final back-substitution and storage of the triangular factor. 
For the ML-CUP setting, where $n$ is fixed and much smaller than $m$, this is quadratic 
in $m$. This is also consistent with the assignment constraint: the algorithm exploits 
the diagonal regularization block instead of treating $\lambda I_m$ as a dense matrix.


% ═════════════════════════════════════════════════════════════
\section{QR with Householder Reflectors}
\label{sec:householder_qr}
% ═════════════════════════════════════════════════════════════

Householder reflectors are the standard stable mechanism for computing QR 
factorizations~\cite{trefethen1997,golub2013}. A Householder reflector is an orthogonal 
matrix of the form

\begin{equation}
\label{eq:householder_def}
    H = I - \tau uu^T,
    \qquad
    \tau = \frac{2}{u^T u},
\end{equation}

where $u\neq 0$. It is symmetric and orthogonal:

\begin{equation}
    H^T = H,
    \qquad
    H^T H = I.
\end{equation}

Given a vector $z\in\R^\ell$, the reflector is chosen so that all components below the 
first one are annihilated. Let

\begin{equation}
    \alpha = -\sign(z_1)\norm{z}_2,
    \qquad
    u = z - \alpha e_1.
\end{equation}

Then

\begin{equation}
    Hz = \alpha e_1.
\end{equation}

Here $\sign(0)$ is taken as $1$. The sign choice avoids cancellation when $z_1$ is 
already close to $\norm{z}_2$, and is therefore important for numerical stability.

\subsection{Dense Householder QR}

For a generic dense matrix $A\in\R^{p\times m}$, Householder QR proceeds column by 
column. At step $k$, one constructs a reflector $H_k$ acting only on rows 
$k,\ldots,p$ such that

\begin{equation}
    H_k A_{k:p,k}
    =
    \begin{bmatrix}
        \alpha_k \\ 0 \\ \vdots \\ 0
    \end{bmatrix}.
\end{equation}

The same reflector is applied to the whole trailing submatrix $A_{k:p,k:m}$. After 
$m$ steps,

\begin{equation}
    H_m\cdots H_2H_1 A =
    \begin{bmatrix}
        R \\ 0
    \end{bmatrix}.
\end{equation}

Since each reflector is orthogonal,

\begin{equation}
    Q^T = H_m\cdots H_1,
    \qquad
    A=QR.
\end{equation}

The important implementation point is that $Q$ is never formed. Only the Householder 
vectors $u_k$ and scalars $\tau_k$ are stored. To compute $Q^T b$, one applies the 
stored reflectors to $b$ in the same order used during the factorization:

\begin{equation}
    b \leftarrow H_k b
    =
    b - \tau_k u_k(u_k^Tb).
\end{equation}

This gives the vector $Q^Tb$, whose first $m$ entries define the right-hand side $c$ 
in Eq.~\eqref{eq:qr_triangular_solve}.

\begin{algorithm}[H]
\caption{Compact Dense Householder QR}
\label{alg:householder_dense}
\begin{algorithmic}[1]
    \Require Matrix $A\in\R^{p\times m}$ with $p\geq m$
    \For{$k=1,\ldots,m$}
        \State $z \gets A_{k:p,k}$
        \State $\alpha \gets -\sign(z_1)\norm{z}_2$
        \State $u \gets z-\alpha e_1$; \quad $\tau \gets 2/(u^Tu)$
        \State $A_{k:p,k:m} \gets A_{k:p,k:m} - \tau u(u^TA_{k:p,k:m})$
        \State Store $(u,\tau)$
    \EndFor
    \State $R \gets \mathrm{triu}(A_{1:m,1:m})$
    \State \Return $R$ and stored reflectors
\end{algorithmic}
\end{algorithm}

Algorithm~\ref{alg:householder_dense} is the mathematically canonical version. It is 
useful for explaining QR, but for the project matrix it should be specialized as 
described next.

\subsection{Structured Householder Row Insertion}
\label{sec:structured_householder}

Using the permuted augmented system~\eqref{eq:permuted_augmented}, the QR factorization 
starts from the already triangular matrix $R_0=\lambda I_m$. For each data row $x_i^T$ 
and scalar right-hand side $y_i$, we maintain an upper triangular matrix $R$ and a 
transformed right-hand side $c$. Inserting one row amounts to reducing

\begin{equation}
    \begin{bmatrix}
        R \\ x_i^T
    \end{bmatrix}
\end{equation}

back to triangular form. This can be done with $m$ Householder reflectors, each with 
support on only two rows: the current triangular row and the active inserted row.

At position $k$, the only entry below the diagonal in column $k$ is the active row 
entry $a_k$. The two-vector

\begin{equation}
    z =
    \begin{bmatrix}
        R_{kk} \\ a_k
    \end{bmatrix}
\end{equation}

is transformed by a two-dimensional Householder reflector

\begin{equation}
    H_k^{(i)} = I_2 - \tau u u^T,
    \qquad
    H_k^{(i)}z =
    \begin{bmatrix}
        \rho \\ 0
    \end{bmatrix}.
\end{equation}

The same reflector is applied to the trailing entries of the two affected rows:

\begin{equation}
    \begin{bmatrix}
        R_{k,k:m} \\
        a_{k:m}
    \end{bmatrix}
    \leftarrow
    H_k^{(i)}
    \begin{bmatrix}
        R_{k,k:m} \\
        a_{k:m}
    \end{bmatrix},
    \qquad
    \begin{bmatrix}
        c_k \\ \eta
    \end{bmatrix}
    \leftarrow
    H_k^{(i)}
    \begin{bmatrix}
        c_k \\ \eta
    \end{bmatrix},
\end{equation}

where $\eta$ is the active right-hand-side entry, initialized as $\eta=y_i$. After the 
$m$ reflectors have been applied, the inserted row has been zeroed and can be discarded. 
The updated $R$ and $c$ are exactly what would be obtained by applying QR to all rows 
processed so far.

\begin{algorithm}[H]
\caption{Structured Householder QR for $[X^T;\lambda I_m]$}
\label{alg:structured_householder}
\begin{algorithmic}[1]
    \Require $X\in\R^{m\times n}$, $y\in\R^n$, $\lambda>0$
    \State $R \gets \lambda I_m$; \quad $c \gets 0\in\R^m$
    \For{$i=1,\ldots,n$}
        \State $a \gets X_{:,i}$; \quad $\eta \gets y_i$
        \For{$k=1,\ldots,m$}
            \State $z \gets [R_{kk},\, a_k]^T$
            \State Build the Householder reflector $H$ such that $Hz=[\rho,\,0]^T$
            \State $[R_{k,k:m};\,a_{k:m}] \gets H [R_{k,k:m};\,a_{k:m}]$
            \State $[c_k;\,\eta] \gets H[c_k;\,\eta]$
            \State Store the reflector parameters if later products with $Q$ or $Q^T$ are required
        \EndFor
    \EndFor
    \State Solve $Rw=c$ by back-substitution
    \State \Return $w$
\end{algorithmic}
\end{algorithm}

\subsection{Complexity and Numerical Properties}

For each inserted row, the inner loop updates trailing row segments of lengths 
$m,m-1,\ldots,1$. Therefore the cost per row is

\begin{equation}
    \sum_{k=1}^m O(m-k+1) = O(m^2),
\end{equation}

and processing all $n$ rows of $X^T$ costs $O(nm^2)$. The final triangular solve costs 
$O(m^2)$. The storage is $O(m^2)$ for the upper triangular factor $R$ and $O(m)$ for the 
right-hand side $c$; if all reflectors are stored in order to apply $Q$ and $Q^T$ later, 
the additional storage is $O(nm)$ because each reflector has two-dimensional support.

The method remains backward stable because it is built entirely from orthogonal 
transformations. In contrast with the normal equations, it does not square the condition 
number. The regularization parameter $\lambda$ also guarantees that the triangular 
factor is nonsingular, since

\begin{equation}
    R^TR = A^TA = XX^T + \lambda^2 I_m \succ 0.
\end{equation}

For very small $\lambda$, the problem is still more ill-conditioned, but the QR method 
is affected through $\kappa(A)$ rather than through $\kappa(A)^2$. This is the main 
reason for using QR as the direct baseline against which the L-BFGS iterations are 
compared.


% ═════════════════════════════════════════════════════════════
\section{Comparison: QR vs.\ L-BFGS}
\label{sec:comparison}
% ═════════════════════════════════════════════════════════════

\textit{(To be completed after A2 implementation.)}


% ═════════════════════════════════════════════════════════════
\appendix
\section{Proofs}
% ═════════════════════════════════════════════════════════════

\begin{lemma}[$\Xhat$ has full column rank]
\label{lemma:fullrank}
If $\lambda > 0$, then $\Xhat^T\Xhat = XX^T + \lambda^2 I_m \succ 0$.
\end{lemma}
\begin{proof}
For any nonzero $v \in \R^m$:
$v^T(XX^T + \lambda^2 I_m)v = \norm{X^T v}^2 + \lambda^2\norm{v}^2 \geq \lambda^2\norm{v}^2 > 0.$
\end{proof}

\begin{lemma}[Eigenvalue shift]
\label{lemma:eigenvalues}
If $\alpha$ is an eigenvalue of $A$, then $\alpha + c$ is an eigenvalue of $A + cI$.
\end{lemma}
\begin{proof}
$Av = \alpha v \iff (A + cI)v = (\alpha + c)v.$
\end{proof}

\begin{lemma}[Singular values of $XX^T$]
\label{lemma:singvals}
The eigenvalues of $XX^T \in \R^{m \times m}$ are 
$\{\sigma_1^2, \ldots, \sigma_n^2, 0, \ldots, 0\}$, where $\sigma_1, \ldots, \sigma_n$ 
are the singular values of $X$. In particular, $\lambda_m(XX^T) = 0$ when $m > n$.
\end{lemma}
\begin{proof}
Let $X = U\Sigma V^T$ be the SVD. Then $XX^T = U\Sigma\Sigma^T U^T$, 
where $\Sigma\Sigma^T = \diag(\sigma_1^2, \ldots, \sigma_n^2, 0, \ldots, 0) \in \R^{m \times m}$.
\end{proof}


% ═════════════════════════════════════════════════════════════
\bibliographystyle{plain}
\begin{thebibliography}{99}

\bibitem{nocedal2006numerical}
J.~Nocedal and S.~J.~Wright.
\newblock \emph{Numerical Optimization}.
\newblock Springer, 2nd edition, 2006.

\bibitem{liu1989limited}
D.~C.~Liu and J.~Nocedal.
\newblock On the limited memory {BFGS} method for large scale optimization.
\newblock \emph{Mathematical Programming}, 45(1):503--528, 1989.

\bibitem{broyden1970}
C.~G.~Broyden.
\newblock The convergence of a class of double-rank minimization algorithms.
\newblock \emph{IMA Journal of Applied Mathematics}, 6(1):76--90, 1970.

\bibitem{hastie2009elements}
T.~Hastie, R.~Tibshirani, and J.~Friedman.
\newblock \emph{The Elements of Statistical Learning}.
\newblock Springer, 2nd edition, 2009.

\bibitem{trefethen1997}
L.~N.~Trefethen and D.~Bau~III.
\newblock \emph{Numerical Linear Algebra}.
\newblock SIAM, 1997.

\bibitem{golub2013}
G.~H.~Golub and C.~F.~Van~Loan.
\newblock \emph{Matrix Computations}.
\newblock Johns Hopkins University Press, 4th edition, 2013.

\bibitem{rodomanov2021}
A.~Rodomanov and Y.~Nesterov.
\newblock New results on superlinear convergence of classical quasi-{N}ewton methods.
\newblock \emph{Journal of Optimization Theory and Applications}, 188(3):744--769, 2021.

\bibitem{jin2022}
Q.~Jin and A.~Mokhtari.
\newblock Non-asymptotic superlinear convergence of standard quasi-{N}ewton methods.
\newblock \emph{Mathematical Programming}, 194(1--2):435--468, 2022.

\bibitem{wolfe1969}
P.~Wolfe.
\newblock Convergence conditions for ascent methods.
\newblock \emph{SIAM Review}, 11(2):226--235, 1969.

\end{thebibliography}

\end{document}
